% Auto-generated from src/content/cv/teaching-statement.mdx — do not edit by hand
\subsection*{Statement of Teaching Philosophy}

**Blake A. Wilson**\par
\medskip

The value of an engineer is fundamentally shifting. In an era where generative AI and automated solvers can rapidly produce code, circuit schematics, and system architectures, the primary role of the engineer is not primarily a designer, but a verifier. My overarching goal as an educator is to train students to be "critical designers" who can rigorously validate automated outputs and understand the physical and computational layers from first principles. To achieve this, my pedagogy relies on testing foundational knowledge in exams, active failure analysis when interacting with generative tools, and inclusive mentorship that treats students as junior colleagues.\par
\medskip

\subsubsection*{The "Verify-First" Classroom}

When students step into the modern engineering landscape, they are handed tools that act as faulty design generators. If we don't teach students how to criticize generative tools and they rely on them too much to retrieve information, they will be ill-equipped to catch the subtle, systemic errors these tools introduce. Therefore, I utilize a backward-design process to build courses where systems orchestration and failure analysis are primary learning outcomes.\par
\medskip

In practice, this means moving beyond traditional rote problem sets. I implement active learning frameworks such as *adversarial justification*. For example, in an algorithms or architecture course, I might provide students with a faulty AI-generated script designed to optimize a data structure or simulate a hardware component. The assignment is not to write the code from scratch, but to act as the senior engineer: they must write the test bench that breaks the AI’s code, mathematically justify why the AI's solution is suboptimal, and provide a correct modification. This forces students to deeply engage with the underlying logic rather than passively trusting a "black box" output.\par
\medskip

\subsubsection*{Scaffolding Mental Models: Tools vs. Truth}

To verify generative tools, students must first possess a rock-solid internal mental model of the fundamentals. I achieve this by distinctly separating how tools are used in preparation versus how mastery is validated:\par
\medskip

\begin{itemize}\setlength{\itemsep}{0.25em}
  \item \textbf{Homework as Orchestration:} I design homework assignments as a training ground for real-world engineering. Here, students are encouraged to use modern tools to build intuition, explore edge cases, and piece together complex workflows. This teaches them \textit{how} to use the tools effectively and safely.
  \item \textbf{Exams as Foundational Validation:} Conversely, I maintain a strict "no-tool" environment for foundational exams. This is not to artificially restrict students, but to accurately assess their individual mastery of core concepts—such as memory management in C or the physics of an electrical circuit. If a student does not possess the foundational truth in their own mind, they cannot act as the "human-in-the-loop" when a tool inevitably hallucinates.
\end{itemize}

\subsubsection*{Mentorship and Inclusive Engineering}

Beyond the lecture hall, my teaching philosophy is heavily informed by my experience leading cross-functional research teams across Oxford, Cambridge, and Harvard. I have formally mentored over 200 researchers, and my approach is always to demystify the "hidden curriculum" of STEM. I focus heavily on creating a welcoming environment for students from non-traditional backgrounds, helping them connect with the engineering community and its principles no matter their background.\par
\medskip

I do this by treating my mentees as junior colleagues from day one. Instead of handing them solved problems, we collaborate on the frontier of what is unknown. We discuss not only the technical mechanisms of our work but the broader societal and ethical landscape of computing. By empowering students to bridge the gap between disciplines and giving them ownership over their work, they transform from students into confident, capable engineers.\par
\medskip

Ultimately, my goal is not just to teach students how to build systems, but to ensure they have the technical rigor and ethical ownership to take professional responsibility for the next generation of engineering.\par
\medskip
